% Generated draft: requires mathematical and source review.
\documentclass[11pt]{article}
\usepackage[margin=25mm]{geometry}
\usepackage{fontspec}
\setmainfont{Latin Modern Roman}
\usepackage{amsmath,amssymb,mathtools,unicode-math}
\setmathfont{Latin Modern Math}
\usepackage{enumitem,needspace,xurl,hyperref}
\setlength{\parindent}{0pt}
\setlength{\parskip}{5pt}
\setlength{\emergencystretch}{3em}
\newcommand{\N}{\mathbb{N}}
\newcommand{\Z}{\mathbb{Z}}
\newcommand{\Q}{\mathbb{Q}}
\newcommand{\R}{\mathbb{R}}
\newcommand{\C}{\mathbb{C}}
\newcommand{\F}{\mathbb{F}}
\newcommand{\A}{\mathbb{A}}
\newcommand{\PP}{\mathbb{P}}
\newcommand{\E}{\mathbb{E}}
\newcommand{\eps}{\varepsilon}
\newcommand{\dd}{\,\mathrm{d}}
\begin{document}
\section*{UnsolvedMath ID: 30001603}
\textbf{Source problem code:} catalogue entry
\subsection{Definitions and mathematical statement}
Let \(X\) be a smooth projective toric variety and let \(E\) be a toric vector bundle on \(X\), meaning a vector bundle equipped with a compatible lift of the torus action that is linear on fibers. The ground field and its characteristic are not specified in the supplied record. With \(\PP(E)\) parametrizing one-dimensional quotients, \(E\) is called nef if the tautological line bundle \(\mathcal O_{\PP(E)}(1)\) has nonnegative degree on every integral projective curve in \(\PP(E)\).

For an integer \(k\geq 1\), use the following jet-spanning convention: \(E\) is \(k\)-jet spanned if, for every closed point \(x\in X\) with maximal ideal \(\mathfrak m_x\subset\mathcal O_{X,x}\), the natural evaluation map
\[
\operatorname{ev}_{x}^{k}\colon H^0(X,E)\longrightarrow E_x/\mathfrak m_x^{k+1}E_x
\]
is surjective, where \(E_x\) denotes the stalk. The intended ground-field and point conventions require confirmation from the source.

\textbf{Question as recorded.} For every smooth projective toric variety \(X\), every nef toric vector bundle \(E\) on \(X\), and every integer \(k\geq1\), does
\[
\tau(E)\geq k
\quad\Longrightarrow\quad
\operatorname{ev}_{x}^{k}\text{ is surjective for every closed point }x\in X
\]
hold? The supplied record also presents the question
\[
\varepsilon(E)\geq k
\quad\Longrightarrow\quad
E\text{ is }k\text{-jet spanned}
\]
as equivalent.

\textbf{Unresolved qualifications.} Neither \(\tau(E)\) nor \(\varepsilon(E)\) is defined in the supplied record. Their normalization, any dependence on a point, and any operation passing from local to global invariants are therefore unspecified. The two displayed implications retain the recorded target, but their numerical hypotheses cannot be made fully precise from the supplied evidence, and the asserted equivalence requires source confirmation.
\subsection{Short English statement}
The recorded conjecture asks whether the bound \(\tau(E)\geq k\) forces a nef toric vector bundle on a smooth projective toric variety to generate every jet of order \(k\) at every point, for each integer \(k\geq1\). The record gives \(\varepsilon(E)\geq k\) as an equivalent hypothesis. Definitions of both numerical invariants and the intended ground field are missing, so this remains a qualified statement-only draft.
\subsection{Sources}
\begin{itemize}
\item Local registry: \texttt{data/unsolvedmath/problems.json}, with \texttt{catalog/problems.json} as the additional-record registry; supplied permanent UnsolvedMath ID \(30001603\), title \emph{Jet Spanning by Nef Toric Vector Bundles}, exact source problem code \texttt{catalogue entry}, in the Oberwolfach Reports Open Problems collection. The supplied record does not identify an individual proposer.
\item Supplied report citation: \emph{Mini-Workshop: Linear Series on Algebraic Varieties (2011)}, Oberwolfach Reports, DOI \texttt{10.4171/owr/2010/45}, pp.\ 2613--2650. Supplied source URL and literature-source link: \url{https://doi.org/10.4171/owr/2010/45}. The registry separately records proposed year \(2010\); the relationship between that year and the citation's \(2011\) designation requires confirmation.
\item Supplied literature assessment, dated exactly \texttt{2026-08-21}: the extracted implication is described as a conjecture, and known toric and line-bundle cases are said not to settle the general vector-bundle formulation. The supplied status is \emph{open}. This is supplied dated curation, not an independently verified literature assessment; its background expressly requests independent verification of time-sensitive status and cited claims.
\item The supplied statement status is \texttt{corrected\_verified}, accompanied by a claim of manual semantic, typographic, and LaTeX review. That registry label does not resolve the missing invariant definitions or constitute independent mathematical verification of this draft.
\end{itemize}
\end{document}
